\chapter{Méthode et résultats}

La méthodologie décrite dans ce chapitre expose l’ensemble des étapes permettant de constituer, d’aligner et
d’exploiter un corpus contenant des textes en français, alsacien, corse et occitan. Le but est d’entraîner
un modèle capable d’identifier automatiquement des
paires de phrases parallèles. Pour ce faire, nous utiliserons des outils d’alignement interlinguistique,
des techniques de représentation sémantique, ainsi que des méthodes statistiques de classification supervisée.

Dans un premier temps, nous utiliserons l’outil \textit{InterText} afin d'aligner les données issues de
l’Atlas des langues régionales du LIMSI. Ce premier travail nous permet d'avoir une autre partie de données pour le corpus,
avec PARCOLAB et la parabole de l'enfant prodigue. Nous calculerons ensuite des scores de similarité textuelle
sur ce corpus.

Les similarités sémantiques issues du modèle SBERT seront
exploitées pour attribuer des scores aux comparaisons interphrastiques sur des données parallèles,
semi-parallèles ou non-parallèles issues du corpus d'entraînement. Ces scores, combinés à d'autres variables
complémentaires mesurant la proximité structurelle et phonétique entre les phrases,
serviront de variables d'entrée pour l'entraînement d'un modèle de régression logistique capable
d'identifier automatiquement les phrases parallèles.

Enfin, nous présenterons une étape complémentaire permettant d’évaluer le bon fonctionnement du modèle en nous
appuyant sur le corpus établi à partir du programme de \citet{tsolakis_2024}. Après validation du modèle sur ce corpus de
référence, il nous faudra étendre le code du modèle afin d'identifier non plus seulement des phrases
parallèles, mais également des fragments semi-parallèles.

\newpage

\section{Alignement manuel des données}

\subsection{InterText}

\href{https://wanthalf.saga.cz/intertext}{InterText} \citep{vondricka_2014} est un environnement conçu pour faciliter
l’annotation et l’alignement de textes multilingues. Développé à l’origine pour les études contrastives et les
sciences du texte, \textit{InterText} permet de comparer des documents appartenant à des langues différentes
et de produire des alignements fiables au niveau phrase ou paragraphe. Pour réaliser ces alignements, \textit{InterText}
se base sur des aligneurs déjà existants
tels que Hunalign \citep{nicolov_2007}, qui exploite des heuristiques basées sur la longueur des phrases et la
fréquence des mots, et TCA2 \citep{hofland_1996, hofland_1998, Vondricka_2010}, un aligneur développé
initialement pour des corpus de langues scandinaves.

\begin{figure}[H]
    \centering
    \includegraphics[width=\textwidth]{ui_intertext.png}
    \caption{L'interface utilisateur d'InterText.}
\label{ui_intertext}
\end{figure}

Son principal avantage réside dans sa capacité à fournir une interface (voir la figure \ref{ui_intertext}) permettant
d’afficher côte à côte deux versions d’un même document (ou de deux documents comparables) afin d’aligner
semi-automatiquement ou manuellement les segments. Concrètement, l'utilisateur
importe deux documents au format XML ou txt, qui sont ensuite segmentés automatiquement en unités textuelles
(phrases, syntagmes ou paragraphes) selon des règles de segmentation paramétrables. Ces segments peuvent
être alignés de trois manières~:

\begin{itemize}
\item \textbf{Alignement automatique}~: l'outil s'appuie sur Hunalign \citep{nicolov_2007} ou TCA2 \citep{hofland_1996, hofland_1998, Vondricka_2010}
pour établir des correspondances probabilistes entre segments sur la base d'indices lexicaux ou structurels.
\item \textbf{Alignement assisté}~: l'utilisateur valide ou corrige les propositions générées
automatiquement. Cette modalité constitue un compromis intéressant entre efficacité et précision,
permettant de corriger les erreurs de l'alignement automatique tout en réduisant la charge de travail manuel.
\item \textbf{Alignement manuel}~: utilisé lorsque la langue ou la variation dialectale rendent
les procédures automatiques peu fiables. C'est cette modalité qui a été principalement mobilisée
dans notre cas, compte tenu de la forte variabilité dialectale des textes de l'Atlas LIMSI.
\end{itemize}

Au terme du processus, \textit{InterText} permet d'exporter ces
alignements sous forme de paires structurées ou de fichiers textes séparés. Pour un corpus 
tel que celui de l’Atlas LIMSI, qui n’a pas été initialement conçu comme un corpus parallèle,
\textit{InterText} représente un outil indispensable permettant de pallier l’absence d’alignement natif.

\subsection{Alignement du corpus LIMSI avec InterText}\label{limsi_intertext}

Le corpus issu de l’Atlas des langues régionales du LIMSI est un corpus dont tous
les textes reposent sur une traduction de la fable
\textit{La bise et le soleil} d'Ésope. Les variations dialectales, les découpages phrastiques
et les divergences d’interprétation rendent l’alignement automatique impossible, ce qui pose
une difficulté supplémentaire à l'alignement de ce corpus.

Dans cette perspective, \textit{InterText} a été mobilisé afin de constituer un corpus parallélisé entre les
variétés régionales. L’alignement s’effectue à partir du texte de référence en français et des versions de
la fable dans les langues régionales de France sur lesquelles nous travaillons.
Pour chaque paire de documents, les segments sont alignés un à un lorsque la correspondance est claire,
ou fusionnés lorsqu'une phrase française correspond à plusieurs segments dans la langue régionale, ou
inversement.

\subsection{Problèmes d'alignement rencontrés}

\begin{figure}[H]
    \centering
    \includegraphics[width=\textwidth]{ui_intertext_errsnttlong.png}
    \caption{L'interface utilisateur d'InterText lors d'une erreur d'alignement (phrase séparée).}
\label{ui_intertext_errsnttlong}
\end{figure}

Durant l'alignement du corpus de l'atlas LIMSI avec InterText, le problème principal rencontré
était dû à des fins de phrases prématurées en langue régionale comparées à la langue française. Dans la phrase 1 de la figure \ref{ui_intertext_errsnttlong}
on peut voir cette erreur, une phrase en français qui se sépare, ici, en 2 phrases en alsacien.
Ces erreurs étaient relativement simples à régler étant donné que la phrase pouvait simplement
être rallongée. En effet, cette erreur ne se présentait que sur des phrases dont le changement de ponctuation ne changeait pas le
sens de la phrase et dont la traduction est restée juste malgré tout.

\begin{figure}[H]
    \centering
    \includegraphics[width=\textwidth]{ui_intertext_errrep.png}
    \caption{L'interface utilisateur d'InterText lors d'une erreur d'alignement (répetition textuelle).}
\label{ui_intertext_errrep}
\end{figure}

Une autre erreur, bien moins grave dans notre cas, était la duplication de texte (voir figure \ref{ui_intertext_errrep}). L'atlas LIMSI
présente parfois lors de son téléchargement des textes qui sont dupliqués dans leur fichier respectif. Ces
erreurs sont très simples à repérer lors de l'alignement et à réparer. En effet, la simple présence
du texte en double est très visible, surtout lorsque celui-ci se duplique du début à la fin après sa
première apparition.

\section{Pré-traitement et analyse du corpus}

\subsection{Séparation du corpus}

Nous commençons par constituer un ensemble de paires de fichiers textes, 
chacune associant un document source en français à son équivalent dans une langue régionale. 
Pour l'Atlas LIMSI, le texte français de référence est systématiquement mis en regard de chacune des versions
dialectales disponibles. Pour la parabole du fils prodigue, le même principe est appliqué. Le corpus ParCoLab,
quant à lui, requiert un
traitement légèrement différent, les fichiers y étant nommés selon une convention incluant l'indication
de langue, le programme identifie dynamiquement les fichiers français et leur associe les fichiers de la
même collection dans les langues cibles.

Chaque paire est ainsi définie par trois éléments~: 
\begin{itemize}
    \item Le fichier source en français
    \item Le fichier cible en langue régionale
    \item Un identifiant permettant de tracer l'origine de la paire au sein du corpus
\end{itemize}
Cette identification est importante, car elle permet, lors des analyses ultérieures, de distinguer les paires
selon leur provenance et donc selon leur profil linguistique et leur niveau de bruit attendu.

Ensuite nous procédons au chargement effectif des fichiers. Chaque fichier est 
lu ligne par ligne, chacune d'entre elles correspondant à une phrase alignée. Le programme vérifie que les deux 
fichiers d'une même paire contiennent le même nombre de phrases, condition impérative pour que l'alignement 
soit considéré comme valide. Si cette condition n'est pas remplie, la paire est écartée. Cette vérification garantit
la cohérence structurelle des données d'entraînement avant toute opération de calcul de score.

\subsection{Création de données parallèles, semi et non-parallèles}

Une fois les paires parallèles chargées, nous prévoyons d'entraîner un modèle supervisé. 
Pour ce faire, il est nécessaire de disposer de
données annotées selon différents degrés de correspondance. Nous distinguons trois catégories~:

\begin{itemize}
    \item Parallèles~:
    \[
    \text{phrase A} \Rightarrow \text{phrase A traduite.}
    \]
    Les deux segments expriment exactement la même information, même si la structure syntaxique varie.
    
    \item Semi-parallèles~:
    \[
    \text{phrase A} \Rightarrow \text{phrase A traduite, mais avec une partie manquante.}
    \]
    Ce cas se rencontre fréquemment lorsque certains détails sont omis ou condensés dans une version dialectale.
    Dans ce mémoire, nous allons devoir créer artificiellement des variations de phrases existantes afin
    de produire des paires ayant un contenu informationnel similaire sans être des traductions exactes.
    Par exemple, une phrase peut être transformée par paraphrasage automatique, par suppression de certains
    éléments d’information (comme la suite d'une phrase après une virgule) ou par reformulation partielle.
    
    \item Non-parallèles~:
    \[
    \text{phrase A} \Rightarrow \text{phrase B distincte sans corrélation sémantique.}
    \]
    Ces paires servent de négatifs dans l’apprentissage supervisé.
\end{itemize}

Le sous-ensemble parallèle est constitué directement à partir des paires chargées précédemment. Si une 
taille maximale a été définie par l'utilisateur, un échantillonnage aléatoire est effectué afin de ne 
retenir qu'un nombre limité de paires. Dans le cas contraire, l'intégralité des paires disponibles est 
conservée. On peut cependant remarquer qu'à l'intérieur du corpus ParCoLab, certains fichiers ne sont pas alignés.
En effet, il y manque des phrases ou des phrases supplémentaires sont présentes.
Les fichiers touchés sont : feltin-20211123\_tr\_gsw.txt (49 phrases FR, 47 phrases LR), feltin-20211123\_tr\_oc.txt (49 FR, 47 LR),
feltin-20220308\_tr\_oc.txt (39 FR, 38 LR), feltin-20220308\_tr\_gsw.txt (39 FR, 38 LR), feltin-20210921\_tr\_gsw.txt (44 FR, 39 LR) et
feltin-20210921\_tr\_oc.txt (44 FR, 39 LR). Par souci de sûreté de l'alignement, nous prenons le nombre de phrases en compte
et décidons de ne pas utiliser les fichiers si le nombre de phrases entre français et langue régionale n'est pas le
même.

Nous créons aussi un sous-ensemble semi-parallèle. La notion de \textit{semi-parallélisme} est importante étant donné
que nous allons travailler sur des corpus comparables qui incluent probablement des segments dont le parallélisme
est suffisant mais trop incertain pour être considéré comme purement parallèle.
Cette notion permettra au modèle de traiter ces cas intermédiaires sans les rejeter systématiquement
comme des paires non-parallèles.
Le sous-ensemble est généré artificiellement à partir des paires parallèles, en 
tronquant le segment en langue régionale, étant donné que nous nous intéressons aux traductions du français
vers ses langues dans notre cas et non l'inverse. Deux stratégies de troncature sont appliquées selon la structure 
de la phrase.
La stratégie privilégiée consiste à identifier la première ponctuation interne de la phrase 
(virgule, point-virgule ou deux-points) et à ne conserver que la portion qui la précède (voir tableau \ref{method1_truncate}). Cette approche est 
motivée par le fait que les transcriptions dialectales de l'Atlas LIMSI présentent fréquemment des omissions 
ou des condensations en milieu ou en fin de phrase, rendant cette troncature linguistiquement plausible.

\begin{table}[H]
\centering
\begin{tabular}{lc}
\hline
\textbf{Phrase} & \textbf{tronquée} \\
\hline
Allura u sù a cuminzzaiu a lusgi, dopu u mumentu, [...]. & non\\
Allura u sù a cuminzzaiu a lusgi, & oui\\
\hline
\end{tabular}
\caption{Tableau présentant la troncature de phrases en langues régionales (méthode 1).}
\label{method1_truncate}
\end{table}

Lorsqu'aucune ponctuation interne n'est disponible, une stratégie de repli consiste à ne conserver que la 
première moitié des mots de la phrase (voir tableau \ref{method2_truncate}). Dans les deux cas, le segment français d'origine est conservé, 
de sorte que la paire résultante présente une correspondance partielle entre les deux langues.

\begin{table}[H]
\centering
\begin{tabular}{lc}
\hline
\textbf{Phrase} & \textbf{tronquée} \\
\hline
Cusci a tramuntana ha ricuniscuiu che u cïu forti de dui ira ben u sù. & non \\
Cusci a tramuntana ha ricuniscuiu che u & oui \\
\hline
\end{tabular}
\caption{Tableau présentant la troncature de phrases en langues régionales (méthode 2).}
\label{method2_truncate}
\end{table}

Le sous-ensemble non-parallèle est quant à lui construit par appariement aléatoire de phrases françaises et 
de phrases en langue régionale provenant de positions différentes dans le corpus. Nous vérifions
systématiquement que les paires ainsi générées ne figurent pas déjà parmi les paires parallèles afin d'éviter 
l'introduction de faux négatifs dans les données d'entraînement. Dans le cas de corpus de petite taille, le 
nombre de combinaisons distinctes disponibles peut être insuffisant pour atteindre la taille cible, ce qui est 
signalé à l'utilisateur.

Les trois sous-ensembles sont enfin fusionnés et mélangés aléatoirement afin de constituer le jeu de 
données complet utilisé pour l'entraînement du classifieur (statistiques du nombre de paires dans le tableau \ref{synthese_ssensembles}).

\begin{table}[H]
\centering
\begin{tabular}{lS}
\hline
\textbf{parallèlisme} & \textbf{Nombre de paires} \\
\hline
Parallèle & 1\ 145 \\
Semi-parallèle & 1\ 145 \\
Non-parallèle & 1\ 145 \\
Total & 3\ 435 \\
\hline
\end{tabular}
\caption{Tableau synthétique du corpus fusionné et mélangé.}
\label{synthese_ssensembles}
\end{table}

\subsection{Calcul des similarités sémantiques avec SBERT}

Une fois l’alignement et la composition de sous-ensemble réalisés, chaque paire de phrases est comparée au moyen d’un modèle reposant sur
l’architecture SBERT \citep{reimers_2019}. SBERT génère pour chaque phrase un embedding, obtenu
par agrégation des sorties de l'encodeur via un mécanisme de \textit{mean pooling}. La similarité entre
deux phrases est ensuite calculée par similarité cosinus entre leurs embeddings respectifs,
produisant un score compris entre 0 et 1.

Comme modèle SBERT, nous décidons d'utiliser le modèle 
\href{https://huggingface.co/sentence-transformers/paraphrase-multilingual-MiniLM-L12-v2}{\textit{paraphrase-multilingual-MiniLM-L12-v2}} \citep{reimers_2019}.
MiniLM est un modèle distillé, il a été entraîné à imiter le comportement d'un modèle plus grand en 
reproduisant ses relations d'attention. Le suffixe L12 indique qu'il possède 12 couches de
transformer, et la dimension de ses embeddings est de 384 dimensions. Un autre modèle nommé 
\href{https://huggingface.co/sentence-transformers/paraphrase-multilingual-mpnet-base-v2}{\textit{paraphrase-multilingual-mpnet-base-v2}} \citep{reimers_2019}
pourrait aussi être intéressant. Les 2 modèles ont été entraînés sur 52 langues~: arabe, bulgare, 
catalan, tchèque, danois, allemand, grec, espagnol, estonien, perse, finlandais, français, français-canadien, galicien, gujarati, hébreu,
hindi, croate, hongrois, arménien, indonésien, italien, japonais, géorgien, coréen, kurde, lituanien, letton, macédonien, mongol, marathi, malaisien, birman,
norvégien bokm{\aa}l, néerlandais, polonais, portugais, portugais-brésilien, roumain/moldave, russe, slovaque, slovène, albanais, serbe, suédois,
thaï, turc, ukrainien, ourdou, vietnamien, mandarin, taïwanais \citep{reimers_2020}.
Le tableau ci-dessous présente les caractéristiques comparatives des deux modèles~:

\begin{table}[H]
\centering
\begin{tabular}{lrr}
\hline
\textbf{Caractéristique} & \textbf{MPNet-base-v2} & \textbf{MiniLM-L12-v2} \\
\hline
Dimension des embeddings & \multicolumn{1}{S}{768} & \multicolumn{1}{S}{384} \\
Nombre de paramètres & $\approx$278M & $\approx$118M \\
Vitesse d'inférence & Modérée & Rapide \\
Qualité sémantique & Supérieure & Inférieure \\
Mémoire nécessaire & Élevée & Faible \\
Langues d'entraînement & \multicolumn{1}{S}{52} & \multicolumn{1}{S}{52} \\
\hline
\end{tabular}
\caption{Tableau comparatif des modèles MPNet et MiniLM.}
\end{table}

Le choix de MiniLM se justifie pour plusieurs raisons complémentaires. Tout d'abord,
les plongements à 384 dimensions restent suffisamment riches pour capturer les similarités sémantiques
entre phrases françaises et régionales. De plus, la légère perte de qualité de MiniLM par rapport à MPNet
aura peu d'impact visible sur les résultats finaux étant donné que le modèle que nous allons entraîner ne se
base pas sur des langues déjà connues et intégrées dans MiniLM ou MPNet. Ce fait vient directement justifier notre
choix, le modèle étant entraîné sur de nouvelles langues (pour l'instant inconnues de celui-ci). Plus de paramètres
et d'entraînement ne feraient que freiner l'entraînement supplémentaire que nous allons appliquer.
Enfin, la réduction de la mémoire nécessaire et l'accélération de l'inférence deviennent des avantages
concrets lors de l'étape de fragmentation, où le nombre de paires à comparer croît significativement.
Les deux modèles étant entraînés sur des données de paraphrase dans plus de 50 langues,
ils sont tous deux capables de projeter des phrases de langues différentes dans un espace sémantique partagé.
C'est indispensable pour comparer directement le français aux langues régionales sans traduction
intermédiaire.

Pour chacune des paires, un score STS est calculé à partir des
plongements SBERT par similarité cosinus.
Un score STS (\textit{Semantic Textual Similarity}) mesure la similarité sémantique entre deux textes
au moyen d'une valeur comprise entre 0 et 1. Plus le score est proche de 1, plus les phrases
expriment la même idée. Plus précisément, le score STS correspond au cosinus
du vecteur de plongement de la première phrase et à celui de la seconde, dans l'espace
vectoriel projeté par le modèle SBERT.
Ces scores, accompagnés des variables complémentaires décrites ci-dessous,
constituent la base d'apprentissage du modèle de régression logistique.

\newpage

\subsection{Principe du modèle de régression logistique}

La régression logistique est un modèle statistique utilisé pour prédire une probabilité appartenant à deux
catégories (dans notre cas~: parallèle / non-parallèle). Cependant, dans notre cas, nous devons prédire 3 catégories.
Nous allons donc créer un modèle de régression logistique multinomiale qui peut prédire les probabilités pour plus
de deux catégories. La régression logistique transforme la sortie en une
probabilité comprise entre 0 et 1 grâce à la fonction sigmoïde~:
\[
\sigma(x) = \frac{1}{1 + e^{-x}}
\]

Dans notre cas, nous allons entraîner un modèle binaire, utilisant la fonction sigmoïde basique, ainsi qu'un modèle
utilisant la fonction softmax de manière à inclure la classe semi-parallèle et à ne pas se limiter aux classes parallèle et
non-parallèle. La fonction softmax se définit ainsi~:
\[
\sigma(\textbf{z})_{i}=\frac{e^{z_{i}}}{\sum_{j=1}^{K}e^{z_{j}}}
\]
Où~:
\begin{itemize}
    \item $\textbf{z} = (z_{1},...,z_{K}) \in \mathbb{R}^{K}$ avec $k > 1$, un tuple de $K$ nombres réels à normaliser en utilisant la fonction $\sigma$
\end{itemize}

La régression logistique constitue un modèle simple pour classifier les paires de phrases selon
qu’elles soient parallèles, semi-parallèles ou non-parallèles (exemples dans le tableau \ref{snt_examples}). Afin d'enrichir la capacité discriminante du
modèle, plusieurs variables d'apprentissage sont combinées en complément des scores STS produits par SBERT.
Les scores issus du modèle SBERT constituent la variable principale. Ils capturent la proximité sémantique
globale entre deux phrases dans un espace vectoriel multilingue partagé.


\begin{table}[H]
\centering
\begin{tabular}{p{0.35\linewidth} p{0.35\linewidth} cS}
\hline
\textbf{Francais} & \textbf{Langue régionale} & \textbf{Classe} & \textbf{Score STS} \\
\hline
Nul ne sera condamné pour des actions ou omissions qui, au moment où elles ont été commises, ne constituaient pas un
acte délictueux d'après le droit national ou international.&
Tombèron d'acòrd qu'aquelh qu'arribariá lo promier a far garar son mantèlh alh viatjairesariá vist coma lo plus fòrt.&
0 & 0,0 \\
On me rétorquera sans doute qu'au moment où la guerre revient en Europe, la France a mieux à faire que de se préoccuper de ses "dialectes".&
Mi riturcerete forse chì, mentre rivene a guerra in Europa, a Francia hà altru à fà chè impenseriscesi di i so « dialetti ».&
1 & 0,8171 \\
Ils sont tombés d'accord que celui qui arriverait le premier à faire ôter son manteau au voyageur serait regardé comme le plus fort.&
Son tombats d'acòrdi que le que arribava le prumièr a far quitar son mantèl al viatjaire&
2 & 0,5371 \\
\hline
\end{tabular}
\smallskip
\small{(0 non-parallèle, 1 parallèle, 2 semi-parallèle)}
\smallskip
\caption{Tableau d'exemples de phrases.}
\label{snt_examples}
\end{table}

\subsection{Score de longueur}

Une variable complémentaire mesure la similarité structurelle entre deux phrases par le biais de leur
longueur respective. Ce score est calculé en divisant la valeur absolue de la différence du nombre de mots
des deux phrases par la longueur de la phrase la plus longue~:
\[ score\ de\ longueur = \frac{\left| len(P_{1})-len(P_{2}) \right|}{max(len(P_{1}),len(P_{2}))} \]
Où $P$ correspond à une phrase.
Un score proche de 0 indique des phrases de longueur similaire, ce qui constitue un indice favorable au
parallélisme. Cette mesure s'inscrit dans la tradition des premières méthodes d'alignement statistique,
comme celle de \citet{gale_1991}, qui exploitaient déjà la corrélation entre longueurs de phrases pour
établir des correspondances probabilistes.

\subsection{Indices de Jaccard} \label{jaccard_idx}

Afin d'exploiter les similarités de surface entre les phrases, notamment les emprunts lexicaux et les
transfuges fréquents entre le français et les langues régionales, des indices de Jaccard
sont calculés sur les bigrammes, trigrammes et quadrigrammes de caractères. L'indice de Jaccard mesure le
chevauchement entre deux ensembles en divisant la taille de leur intersection par la taille de leur union.
Soit~:
\[ J(A,B)=\frac{A \cap B}{A \cup B}=\frac{|A \cap B|}{|A| + |B| - |A \cap  B|} \]
où $A$ et $B$ sont des ensembles dont l'indice mesure le chevauchement.
Avant ce calcul, les phrases sont normalisées par suppression de la casse, de la ponctuation et des accents,
afin de réduire le bruit lié aux variations orthographiques, très fréquentes dans les langues
régionales dont la normalisation est parfois incomplète.

\subsection{Représentations phonétiques}

Une dernière variable exploite les similarités phonétiques entre les phrases via la bibliothèque Python
jellyfish \citep{turk_2013}. Les phrases sont d'abord encodées selon l'algorithme Metaphone \citep{philips_1990},
qui représente les mots selon leur prononciation approximative de sorte que des mots phonétiquement proches
partagent une même représentation. Par exemple, les mots français «~mer~» et «~mère~» produiraient des
représentations Metaphone \citep{philips_1990} similaires (voir tableau \ref{metaphone}). 

\begin{table}[H]
\centering
\begin{tabular}{lcS}
\hline
\textbf{Mots} & \textbf{Représentations Metaphone} & \textbf{Score} \\
\hline
mer / mère & MR / MR & 1,00 \\
Chant / Champ & XNT / XMP & 0,33 \\
Chant / Blé & XNT / BL & 0,00 \\
\hline
\end{tabular}
\caption{Tableau d'exemples de scores Metaphone \citep{philips_1990}.}
\label{metaphone}
\end{table}

Un indice de Jaccard est ensuite calculé sur ces représentations
phonétiques. Cet indice n'est pas calculé sur les n-grammes comme les indices de Jaccard présentés précédemment mais
réalise une simple comparaison ensembliste entre les ensembles de code Metaphone comme défini mathématiquement dans \ref{jaccard_idx}.
Cette variable est particulièrement pertinente dans notre contexte, car les langues régionales de
France partagent souvent avec le français des racines étymologiques communes dont la proximité est davantage
perceptible à l'oral qu'à l'écrit, notamment en occitan ou en corse. Elle permet ainsi de capturer des
correspondances que les mesures purement lexicales ou sémantiques pourraient manquer.

Le modèle apprend ainsi à combiner ces différentes variables pour reconnaître les configurations
typiques des phrases parallèles et à distinguer les paires non-parallèles. L'entraînement sur des données
semi-parallèles devrait rendre le modèle capable d'identifier ce type de segments, dont la présence
est particulièrement fréquente dans les corpus dialectaux comme celui de l'Atlas LIMSI.

\subsection{Difficultés - phonétique}

L'utilisation de la bibliothèque jellyfish pour le calcul des représentations phonétiques a soulevé plusieurs difficultés pratiques.
L'algorithme Metaphone sur lequel elle s'appuie a été conçu pour l'anglais et encode les mots selon des règles phonétiques propres à cette langue.
Son application directe à des textes en français et en langues régionales produit donc des représentations approximatives, parfois
erronées. Par exemple,
le digramme «~ch~» en français ne correspond pas phonétiquement aux conventions Metaphone qui lui attribuent /K/ phonétiquement.
De même, les caractères accentués fréquents dans les langues romanes — tels que «~è~», «~à~» ou «~ï~» — ne sont pas gérés nativement par
l'algorithme. L'étape de normalisation Unicode préalable, qui supprime les diacritiques par décomposition, atténue partiellement ce problème, mais
introduit en contrepartie une perte d'information phonétique. Une alternative plus rigoureuse aurait consisté à utiliser des transcriptions
phonétiques basées sur l'API (alphabet phonétique international), mais aucun outil libre permettant de l'utiliser n'existe à notre connaissance.
Malgré ces limites, l'indice de Jaccard phonétique apporte une contribution
positive au modèle, ce qui suggère que même une représentation phonétique approximative capture un signal utile pour les langues partageant des
racines étymologiques communes avec le français.

\section{Apprentissage supervisé}

\subsection{Application programmatique}

Avant tout calcul, une étape de normalisation est appliquée aux textes. Celle-ci consiste à convertir les 
caractères en minuscules, à supprimer les accents par décomposition Unicode, puis à retirer la ponctuation. 
Cette normalisation est particulièrement importante dans le contexte des langues que nous étudions.
Celles-ci possèdent, en effet, des variations orthographiques dépendant du dialecte qui doivent être normalisées
sous peine de poser problème plus tard.

Enfin, une vérification de cohérence est effectuée sur la première paire chargée (voir la figure \ref{sanity_check}), 
affichant les valeurs de chacun de ces indicateurs afin de s'assurer du bon fonctionnement des fonctions 
avant le traitement de l'ensemble du corpus.

\begin{figure}[H]
    \centering
    \includegraphics[width=\textwidth]{sanity_check.png}
    \caption{Vérification de cohérence sur la première paire de phrases.}
    \label{sanity_check}
\end{figure}

Une fois les variables de surface calculées, nous calculons un score de similarité
sémantique à l'aide du modèle SBERT. Le modèle est d'abord chargé en mémoire, puis appliqué à l'ensemble
des paires constituant le jeu de données normalisées.

Pour chaque paire de phrases, le modèle génère un plongement vectoriel pour chacun des deux segments,
c'est-à-dire une représentation numérique de la phrase dans un espace sémantique multilingue.
La similarité entre les deux phrases est ensuite calculée par similarité cosinus entre leurs plongements
respectifs, produisant un score compris entre 0 et 1. Pour optimiser l'utilisation de la mémoire,
le calcul est effectué par lots plutôt que par phrase. Les scores légèrement négatifs, qui peuvent
ponctuellement résulter du calcul de la similarité cosinus, sont ramenés à zéro afin de maintenir une
échelle de valeurs cohérente.

Ces scores constituent la variable principale du modèle de régression logistique, dans la mesure où ils
capturent la proximité sémantique globale entre les phrases françaises et leurs équivalents en langue
régionale dans un espace vectoriel commun, sans nécessiter de traduction intermédiaire entre les deux
langues.

\begin{table}[H]
\centering
\begin{tabular}{l p{0.35\linewidth} p{0.35\linewidth}}
\hline
\textbf{Indice} & \textbf{FR} & \textbf{LR} \\
\hline
0 &
Alors, la bise s'est mise à souffler de toute sa force mais plus elle soufflait, plus le voyageur serrait son manteau autour de lui et à la fin, la bise a renoncé à le lui faire ôter.&
Cusci a tramuntana ha ricuniscuiu che u cïu forti de dui ira ben u sù. \\
1 &
Ils sont tombés d'accord que celui qui arriverait le premier à faire ôter son manteau au voyageur serait regardé comme le plus fort.&
Tombèron d'accòrdi qu'aquel qu'arribariá le prumièr a far quitar \\
2 &
Dans la tradition républicaine, en effet, il est censé n’y avoir chez nous qu’un seul peuple, le peuple français, dont l’unité a valeur constitutionnelle.&
In a tradizione republicana, di fatti, ci saria ind’è noi un populu solu, u populu francese, di u quale l’unità hà valore custituziunale. \\
\hline
\end{tabular}
\caption{Phrases et leur indice pour comprendre le tableau \ref{variables_extract}.}
\label{jaccard_extract}
\end{table}

\begin{table}[H]
\centering
\begin{tabular}{lcSSSSSS}
\hline
\textbf{Indice} & \textbf{Classe} & \textbf{SBERT} & \textbf{longueur}  & \textbf{jaccard 2} & \textbf{jaccard 3} & \textbf{jaccard 4} & \textbf{jaccard p} \\
\hline
0 & 0 & 0,4010 & 0,5946 & 0,2190 & 0,0435 & 0,0094 & 0,0910 \\
1 & 2 & 0,4257 & 0,6087 & 0,3790 & 0,1901 & 0,1195 & 0,2083 \\
2 & 1 & 0,8039 & 0,0416 & 0,4286 & 0,1822 & 0,0958 & 0,2000 \\
\hline
\end{tabular}
\\
\smallskip
\small{(classes -- 0~: non-parallèle, 1~: parallèle, 2~: semi-parallèle)}
\smallskip
\caption{Extrait du tableau des variables}
\label{variables_extract}
\end{table}

L'ensemble des scores calculés précédemment est ensuite rassemblé dans un tableau structuré, où chaque
ligne correspond à une paire de phrases et chaque colonne à une variable descriptive. Pour chaque paire,
le tableau enregistre les phrases source et cible, l'identifiant de la paire de langues, la catégorie de
parallélisme sous forme textuelle, ainsi qu'une étiquette binaire opposant les paires parallèles et
semi-parallèles d'un côté aux paires non-parallèles de l'autre. Cette étiquette binaire constitue la
variable cible du modèle de régression logistique. À ces informations s'ajoutent l'ensemble des variables
d'apprentissage calculées précédemment~: le score de similarité sémantique SBERT, le score de longueur,
les indices de Jaccard sur les bigrammes, trigrammes et quadrigrammes de caractères, ainsi que l'indice de
Jaccard phonétique (ressemblant au tableau \ref{variables_extract}). Ce tableau constitue la base d'apprentissage complète qui sera transmise au classifieur
lors de l'étape d'entraînement.

\subsection{Analyse des données par classe}

\begin{figure}[H]
    \centering
    \includegraphics[width=\textwidth]{feature_distributions_kde.png}
    \caption{Graphe KDE de distribution des caractéristiques par classe.}
    \label{feature_distributions_kde}
\end{figure}

Ces graphiques (figure \ref{feature_distributions_kde}) présentent les distributions de densité des six variables
d'apprentissage selon les trois catégories de parallélisme, sous forme de courbes d'estimation par noyau (\textit{kernel density estimation} ou KDE).
Les tirets verticaux indiquent la médiane de chaque classe.

Le score SBERT constitue la variable dont le pouvoir discriminant est le plus visible. Les trois distributions
présentent des profils clairement distincts. La classe non-parallèle est fortement concentrée dans les
valeurs basses avec un pic prononcé autour de 0,10 à 0,15, et une médiane proche de 0,18. La classe
semi-parallèle présente une distribution bimodale avec une augmentation autour de 0,15 et une autre autour
de 0,38, reflétant l'hétérogénéité des paires tronquées. Certaines troncatures importantes produisent des
scores faibles proches de ceux des paires non-parallèles, tandis que les troncatures légères conservent une
proximité sémantique élevée. La classe parallèle se distingue par une distribution étalée vers les valeurs
hautes avec un mode principal autour de 0,50 et une médiane nettement supérieure aux deux autres
classes à 0,60. La superposition partielle entre les trois distributions, particulièrement
entre les classes semi-parallèle et non-parallèle dans la région 0,10 à 0,35, explique une partie des
erreurs de classification observées dans les matrices de confusion (voir figures \ref{confusion_binary} et
\ref{confusion_multiclass} plus bas).

Le score de longueur présente un comportement atypique par rapport aux autres variables. Les classes semi et non-parallèles
ont une médiane autour de 0,57 à 0,59 avec une concentration aux alentours de 0,5 et 0,9 respectivement.
À l'inverse, les phrases parallèle ont un score beaucoup plus bas.
Ce score montre que plus la différence de longueur entre la phrase source et la phrase cible est basse, plus la phrase
semble être parallèle. Ce score est utile pour remarquer le parallélisme d'une phrase assez rapidement. En effet,
si le score de longueur est bas entre une phrase en langue source et sa contrepartie en langue cible, il y a de fortes
chances que celles-ci soient parallèles. Cette variable pourrait être utilisée comme un premier filtre simple pour extraire
les phrases potentiellement parallèles des autres. Cependant il est à noter que ce score seul ne prouve pas un
parallélisme. Si nous nous basions uniquement sur la différence de longueur entre les phrases, alors nous
pourrions très facilement tomber sur de faux positifs dans le cas de textes comme des poèmes où la longueur des phrases
peut, parfois, dépendre d'un choix stylistique.

Le score de Jaccard sur les bigrammes montre une meilleure séparation que les variables de trigrammes et de quadrigrammes. La
classe non-parallèle est concentrée entre 0 et 0,40 avec un pic autour de 0,26, tandis que la classe parallèle
présente une distribution nettement plus étalée avec un second mode visible autour de 0,40 à 0,42,
traduisant les chevauchements lexicaux fréquents entre le français et ses langues régionales.
La classe non-parallèle occupe une position bimodale qui se joint assez bien à la courbe de la classe
semi-parallèle. Cette proximité entre les deux courbes est intéressante car elle présente une superposition
des deux classes dans le cas des scores de Jaccard sur les bigrammes. Cette superposition peut provenir de faux négatifs
dans la classe semi-parallèle qui ont été considérés comme des phrases non-parallèles.

Le score de Jaccard pour les trigrammes, quadrigrammes et phonétique présente des distributions très
similaires entre elles. Les trois classes sont massivement concentrées dans les valeurs très
basses, entre 0 et 0,05 à 0,10, avec des pics de densité élevés en 0, ce qui explique les
valeurs d'axe Y atteignant respectivement 10,2, 5,8 et 2 en densité. Cette concentration reflète
le fait que les langues régionales, en particulier l'alsacien, partagent relativement peu de séquences
de caractères longues avec le français malgré leurs racines communes (pour voir cette différence plus en détail,
voir les figures \ref{feature_distributions_kde_roman} et \ref{feature_distributions_kde_gsw}).
La classe parallèle se distingue
par une queue de distribution légèrement plus étendue vers les valeurs positives et une médiane légèrement
supérieure, visible sur les tirets verts, mais la séparation entre classes reste faible pour ces trois
variables prises isolément. Leur combinaison avec le score SBERT au sein du modèle de régression
logistique, leur confère une utilité discriminante accrue, notamment pour les paires où le score SBERT
ou les scores Jaccard ne sont pas suffisants pour trancher sur le parallélisme entre deux phrases.

\begin{figure}[H]
    \centering
    \includegraphics[width=\textwidth]{fig_kde_romance_classes.png}
    \caption{Graphe KDE de distribution des caractéristiques par classe (langues romanes).}
    \label{feature_distributions_kde_roman}
\end{figure}

On peut remarquer que la distribution des caractéristiques par classe est très similaire pour les langues
romanes (figure \ref{feature_distributions_kde_roman}) comparés aux distributions générales. De plus nous pouvons observer que les pics, souvent présents
aux mêmes endroits, sont plus prononcés que sur les distributions générales.

\begin{figure}[H]
    \centering
    \includegraphics[width=\textwidth]{fig_kde_gsw_classes.png}
    \caption{Graphe KDE de distribution des caractéristiques par classe (langues germaniques).}
    \label{feature_distributions_kde_gsw}
\end{figure}

L'alsacien (seule langue germanique du jeu de données) présente, lui, de vastes différences par rapport aux graphes généraux.
On remarque par exemple que les indices de Jaccard ne sont pas, ou presque, distribués sur leurs graphes. Cette distribution
de l'indice très pauvre prouve le manque de séquences de caractères longues entre le français et l'alsacien, les bigrammes étant
plus distribués que les trigrammes, eux-mêmes plus distribués que les quadrigrammes. Ceux-ci restent
plutôt dans les valeurs basses et rapprochés entre eux. De même les scores SBERT sont très rapprochés et les paires
parallèles ont un score SBERT beaucoup plus bas avec une moyenne à 0,41 contre 0,63 pour les langues romanes.
La distribution du score SBERT pour la classe parallèle est d'ailleurs très similaire à
celle de la classe semi-parallèle. Nous pouvons en conclure que les deux classes sont plus floues entre elles dans le
cas des langues germaniques que dans le cas des langues romanes. La variable phonétique présente elle aussi des valeurs
basses pour la plupart des paires. Ces valeurs montrent la vaste différence phonétique entre le français et l'alsacien.

\begin{figure}[H]
    \centering
    \includegraphics[width=\textwidth]{pie_by_language.png}
    \caption{Distribution des classes par paires de langues.}
    \label{pie_by_language}
\end{figure}

Les diagrammes circulaires (figure \ref{pie_by_language}) présentent la distribution des classes au sein du jeu de données
d'entraînement, agrégée par paire de langues. Les effectifs indiqués entre parenthèses (fr-co~: n~=~1\,186,
fr-gsw~: n~=~834, fr-oc~: n~=~1\,415) confirment le déséquilibre déjà observé dans l'analyse quantitative du
corpus~: l'occitan est la langue la mieux représentée, suivi du corse puis de l'alsacien. Ce déséquilibre
est directement imputable à la surreprésentation des variétés occitanes dans l'Atlas LIMSI, dont les 47
textes dialectaux constituent la majorité des paires chargées. Plus largement, cette majorité peut aussi être
imputée par le fait que l'occitan est la langue régionale la plus représentée en France, celle-ci s'étendant
sur la majorité du sud du pays.

Pour les trois paires de langues, la
distribution entre les trois classes est proche de l'équilibre, ce qui est cohérent
avec la stratégie de construction du jeu de données décrite dans la méthodologie, où les sous-ensembles
parallèle, semi-parallèle et non-parallèle sont construits à partir du même ensemble de paires parallèles
de base. On observe néanmoins de légères variations. Pour la paire fr-gsw, la classe non-parallèle est
légèrement surreprésentée à 35,5~\% au détriment de la classe parallèle à 32,3~\%, ce qui peut s'expliquer
par les contraintes de génération des paires non-parallèles dans de petits corpus. En raison du nombre limité
de phrases disponibles, le programme génère plus facilement des paires négatives distinctes que dans
des corpus plus volumineux. Pour la paire fr-oc, la distribution est quasi parfaitement équilibrée à
34,3~\% pour les classes parallèle et semi-parallèle, et 31,4~\% pour la classe non-parallèle, ce qui
s'explique par la taille plus importante de ce corpus permettant un échantillonnage plus équilibré.
Ces distributions globalement équilibrées sont bonnes pour l'entraînement du
classifieur, évitant les biais dus à des déséquilibres importants entre classes.

\begin{figure}[!h]
    \centering
    \includegraphics[width=\textwidth]{sbert_by_language_pair.png}
    \caption{Score SBERT médian par paires de langues.}
    \label{sbert_by_language_pair}
\end{figure}

Ce graphique (\ref{sbert_by_language_pair}) en barres horizontales présente les scores SBERT médians par paire de langues et par catégorie
de parallélisme.

La hiérarchie attendue entre les trois classes est respectée pour l'ensemble des paires de langues~: les
paires parallèles obtiennent systématiquement les scores les plus élevés, suivies des paires semi-parallèles,
puis des paires non-parallèles. Ce résultat semble valider le fait que le modèle SBERT capture efficacement les
différences de parallélisme entre catégories indépendamment de la langue régionale cible.

La paire fr-co présente les scores les plus élevés pour les classes parallèle et semi-parallèle, atteignant
respectivement 0,69 et 0,49.
Cela confirme que le corse, de par sa proximité étymologique avec le
français et sa structure romane, génère des paires
parallèles dont la similarité sémantique vectorielle est plus élevée que pour les autres langues. La paire
fr-oc arrive en deuxième position avec un score parallèle de 0,63 et un score semi-parallèle
de 0,45.
La paire fr-gsw présente les scores les plus bas pour la classe parallèle (0,49), ce qui
reflète la distance structurelle entre le français et l'alsacien.

Les scores non-parallèles sont en revanche remarquablement stables entre les trois paires de langues,
oscillant autour de 0,18 à 0,19. Cette stabilité est attendue, étant donné que les paires non-parallèles, construites
par appariement aléatoire de phrases, ne présentent aucun lien sémantique systématique quelle que soit
la langue régionale impliquée. L'écart entre le score parallèle et le score non-parallèle, qui constitue
la marge de discrimination effective du modèle, est plus favorable pour le corse (0,51) que pour
l'occitan (0,45) ou l'alsacien (0,30), ce dernier présentant la classification la plus difficile.
Cela est sûrement dû au manque de données en alsacien présentes dans le corpus d'entraînement comparé aux
deux autres langues. Cela pourrait aussi être dû à la différence entre l'alsacien et les autres langues
régionales, étant donné que celui-ci est une langue à racine germanique et non romane.

\begin{figure}[H]
    \centering
    \includegraphics[width=\textwidth]{sbert_distribution.png}
    \caption{Distribution des scores SBERT par classe (parallèles, semi et non-parallèles).}
    \label{sbert_distribution}
\end{figure}

Ce double graphique (\ref{sbert_distribution}) présente la distribution des scores SBERT selon les trois catégories étudiées.
L'histogramme de gauche révèle une structure cohérente avec les attentes théoriques. Les paires
non-parallèles se concentrent massivement dans les valeurs basses, entre 0 et 0,3, avec un pic prononcé
autour de 0,05 à 0,10, reflétant l'absence de lien sémantique entre des phrases appariées aléatoirement.
Les paires semi-parallèles présentent une distribution plus étalée, centrée autour de 0,3 à 0,5,
correspondant à la conservation d'une proximité sémantique avec
le segment français mais avec moins de contenu informationnel. Les paires parallèles et semi-parallèles se distinguent nettement
avec une concentration importante dans les valeurs hautes, proche de 1,0, confirmant que le score SBERT
capture efficacement la concordance sémantique entre des traductions proches ou exactes même dans un contexte multilingue
régional. Le graphe de boîtes à moustaches représente les données pour 3\,435 paires de phrases
(soit 1\,145 par classe). Ce graphe synthétise ces observations, la médiane des paires
non-parallèles se situe autour de 0,17, celle des paires semi-parallèles autour de 0,40, et celle des
paires parallèles autour de 0,60. On notera la présence de nombreux points aberrants au-dessus de 0,55 pour
la classe non-parallèle, correspondant vraisemblablement à des phrases thématiquement proches appariées
par hasard lors de la construction du sous-ensemble négatif. La variance importante de la classe parallèle,
visible dans l'amplitude de la boîte,
reflète quant à elle l'hétérogénéité dialectale du corpus, les variétés LIMSI présentant des scores nettement
inférieurs aux textes ParCoLab.

\subsection{Entraînement du modèle de régression logistique}

Une fois le jeu de données constitué et les variables d'apprentissage calculées, il devient possible
d'entraîner un modèle de régression logistique capable d'identifier automatiquement les paires de phrases parallèles ou
semi-parallèles au sein d'un corpus comparable non aligné. Avant l'entraînement proprement dit, une phase exploratoire
permet de vérifier la cohérence du jeu de données.
Les statistiques descriptives calculées par classe confirment que les variables discriminent bien les trois
catégories. Les paires parallèles obtiennent systématiquement les scores SBERT les plus élevés, suivies
des paires semi-parallèles, puis des paires non-parallèles. Cette hiérarchie, observable également dans
les graphiques produits lors de la constitution du corpus, valident l'utilisation des variables retenues pour
l'apprentissage.

Les données sont ensuite divisées en un ensemble d'entraînement et un ensemble de test selon un découpage
stratifié de 80/20, garantissant que les trois classes sont représentées de manière proportionnelle dans
chacun des deux ensembles. Les variables sont normalisées par standardisation, une étape indispensable pour
la régression logistique. Nous entraînons un scaler
ajusté sur les données d'entraînement que nous sauvegardons afin de pouvoir être appliqué de manière cohérente lors
de l'inférence sur de nouvelles données.

Deux modèles distincts sont entraînés. Le premier est un modèle binaire qui oppose les paires parallèles
et semi-parallèles aux paires non-parallèles. Ce modèle constitue le c{\oe}ur opérationnel du système. Appliqué
à un corpus comparable, il permet de filtrer les paires suffisamment parallèles pour
être exploitables. Le second est un modèle multiclasse, qui distingue les paires
parallèles, semi-parallèles et non-parallèles. Ce modèle plus fin permet, pour les paires retenues par
le modèle binaire, de distinguer les correspondances exactes des correspondances partielles. Cette information
peut être utile pour orienter leur usage ultérieur dans l'entraînement d'autres systèmes. Les
deux modèles utilisent une pondération équilibrée des classes afin de pallier d'éventuels déséquilibres
dans le jeu de données, et sont évalués par validation croisée à cinq plis sur les données
d'entraînement avant leur application sur le jeu de test.

Une analyse complémentaire ventile les performances des modèles de régression logistique par langue régionale, en distinguant
le corse, l'alsacien et l'occitan. Cette séparation est importante dans notre contexte, car les trois
langues présentent des profils très différents en termes de proximité avec le français.

Enfin, une fonction d'inférence calcule l'ensemble des variables pour une paire donnée, les normalise à l'aide du scaler
sauvegardé, puis retourne la prédiction du modèle accompagnée d'un score de confiance. Elle peut être
appliquée à l'ensemble des combinaisons de phrases d'un corpus comparable non aligné, produisant une
liste de paires candidates classées selon leur parallélisme. Cette étape constitue le point
d'articulation entre le travail de classification réalisé sur les données d'entraînement
et l'application opérationnelle du modèle sur le corpus Wikipédia utilisé pour la validation.

\subsection{Difficultés - entraînement}

La sauvegarde des modèles entraînés et du scaler de normalisation a été réalisée
à l'aide de la bibliothèque joblib. Cette approche, bien que courante et documentée dans
l'écosystème scikit-learn, soulève plusieurs difficultés. La première difficulté est liée à la
dépendance stricte aux versions des bibliothèques utilisées lors de
l'entraînement. Un fichier \textit{.joblib} produit avec une version donnée de
scikit-learn n'est pas garanti d'être désérialisable avec une autre version de la librairie.
Cela tient au fait que joblib sérialise non seulement les
paramètres appris du modèle mais également la structure interne des objets
Python tels qu'ils existent dans la version de la bibliothèque au moment de la
sauvegarde. Si cette structure change entre deux versions, la désérialisation
échoue ou produit un objet incohérent. Dans le cadre de ce mémoire,
cette contrainte implique que le déploiement du modèle sur une autre machine
nécessite de reproduire exactement l'environnement logiciel utilisé lors de
l'entraînement, ce qui peut s'avérer complexe sans gestion explicite des
dépendances ou un environnement figé. La
deuxième difficulté est d'ordre sécuritaire. La désérialisation d'un fichier \textit{.joblib}
repose sur le module pickle de Python, qui est notoirement
vulnérable à l'exécution de code arbitraire. Le processus de désérialisation
reconstruit des objets en exécutant du code embarqué dans le fichier
binaire, sans contrôle préalable sur la nature de ce code. Bien que cette
vulnérabilité soit sans conséquence dans un contexte de recherche où les
fichiers sont produits et consommés par le même environnement, elle constitue
une limite importante si le modèle venait à être distribué ou intégré dans une
chaîne de traitement partagée. La troisième difficulté concerne le couplage entre le modèle et le scaler.
Dans notre pipeline, le
modèle de régression logistique et le scaler de standardisation sont sauvegardés
séparément dans deux fichiers \textit{.joblib} distincts. Cette séparation impose que les
deux fichiers soient toujours chargés conjointement et dans le bon ordre lors de
l'inférence. Le scaler doit être appliqué en premier sur les variables brutes avant que
celles-ci ne soient transmises au classifieur. Si l'un des deux fichiers est
absent, corrompu ou remplacé par une version incompatible (dans le cas où le scaler se base sur un jeu de données différent, par exemple)
les prédictions produites deviennent erronées, sans qu'aucune erreur ne soit relevée. Ce type de défaillance
silencieuse est particulièrement difficile à détecter en production. Une approche plus robuste consisterait à encapsuler le scaler et le modèle
dans un pipeline scikit-learn unique via \textit{sklearn.pipeline.Pipeline}, qui garantit que la normalisation est toujours appliquée
avant la classification et permet de sauvegarder l'ensemble de manière plus stable.


\subsection{Analyse des caractéristiques}

\begin{figure}[H]
    \centering
    \includegraphics[width=0.85\textwidth]{feature_correlation.png}
    \caption{Corrélation des variables utilisé par le modèle.}
    \label{feature_correlation}
\end{figure}

La matrice de corrélation (\ref{feature_correlation}) entre variables d'apprentissage révèle plusieurs relations
importantes. La corrélation négative entre le score SBERT et le score de longueur (-0,58) est attendue,
une asymétrie de longueur importante entre les deux phrases tend à signaler une correspondance
imparfaite, ce qui se traduit par un score sémantique plus faible. Entre les indices de Jaccard,
on observe une corrélation très élevée entre les trigrammes et les quadrigrammes (0,97), ce qui indique une
forte redondance entre ces deux variables. En effet, ces deux mesures capturent des chevauchements de surface
et apportent une information quasi identique au modèle. La corrélation entre
l'indice de Jaccard phonétique et les trigrammes (0,76) et quadrigrammes (0,71) est également notable,
suggérant que les similarités phonétiques et les similarités de surface en caractères covarient fortement
dans notre corpus. Cela est probablement dû au fait que les langues régionales romanes partagent avec le français des
séquences de caractères proches de leur représentation phonétique. Le score SBERT présente des corrélations
modérées et positives avec les indices de Jaccard (entre 0,35 et 0,47), indiquant que les deux familles de
variables mesurent des aspects complémentaires mais seulement partiellement liés. Cette
structure de corrélation justifie le choix de conserver l'ensemble des variables malgré les
redondances partielles, la régression logistique étant capable de pondérer leur contribution respective.

\begin{figure}[H]
    \centering
    \includegraphics[width=\textwidth]{feature_importance.png}
    \caption{Importance des variables utilisées par le modèle binaire.}
    \label{feature_importance}
\end{figure}

Le diagramme \ref{feature_importance} présente les coefficients des caractéristiques pour le modèle binaire.
Chaque coefficient $w_{j}$ d'une caractéristique $x_{j}$ représente l'impact de celle-ci sur le 
logarithme du rapport de cotes (\textit{odds-ratio}) de la classe cible. Plus précisément, pour le modèle binaire, la prédiction s'effectue selon :
\[
\log\left(\frac{P(y=1 \mid \mathbf{x})}{P(y=0 \mid \mathbf{x})}\right) = w_{0} + w_{1} x_{1} + \cdots + w_{n} x_{n}
\]
Un coefficient positif signifie qu'une augmentation de la caractéristique augmente la probabilité d'appartenir à la classe positive (parallèle/semi-parallèle).
À l'inverse, un coefficient négatif indique que la caractéristique est associée à la classe négative (non-parallèle).

L'indice de Jaccard sur les quadrigrammes de caractères présente de loin le coefficient le
plus élevé (environ 3,0), dépassant significativement le score SBERT (environ 1,4) et l'indice phonétique
(environ 1,0). Ce résultat suggère que le chevauchement de séquences de quatre caractères consécutifs
constitue le signal le plus discriminant pour identifier les paires parallèles dans notre corpus, davantage
encore que la similarité sémantique. Cela s'explique par la nature des langues étudiées~:
l'occitan, le corse et l'alsacien partagent avec le français de nombreuses racines lexicales dont la graphie
est suffisamment proche pour que les quadrigrammes constituent un indicateur de
correspondance fort. Le score SBERT arrive en deuxième position, confirmant son rôle central mais montrant qu'il
n'est pas le seul déterminant fiable dans notre contexte. L'indice phonétique contribue positivement, mais dans
une moindre mesure, tandis que les bigrammes de caractères ont un effet positif faible. Le coefficient négatif
des trigrammes (environ -0,4) est en apparence paradoxal mais s'explique par la forte similarité avec les
quadrigrammes observée dans la matrice de corrélation. Le modèle compense partiellement la redondance entre
ces deux variables en attribuant un poids négatif à l'une d'elles.
Enfin, le score de longueur présente un coefficient très proche de zéro, indiquant qu'une fois les autres variables contrôlées, l'asymétrie de
longueur n'apporte que peu à la classification.

\begin{figure}[H]
    \centering
    \includegraphics[width=\textwidth]{feature_importance_binary_or.png}
    \caption{Rapports de cotes des caractéristiques dans le modèle binaire.}
    \label{feature_importance_or}
\end{figure}

Le rapport de cotes $e^{w_{j}}$ permet
d'interpréter l'effet de la caractéristique dans l'espace probabiliste. Par exemple, un rapport de 3 signifie qu'une augmentation d'un écart-type de la
caractéristique multiplie les chances d'être dans une classe par 3. Une caractéristique éloignant d'une classe d'un rapport inférieur à 1
(en rouge sur le graphe \ref{feature_importance_or}).
Ainsi, le rapport de cotes est plus intuitif que le coefficient brut pour
communiquer le poids qu'une caractéristique prend au sein du modèle.

Le diagramme \ref{feature_importance_or} présente les rapports de cotes pour les caractéristiques
au sein du modèle binaire. On remarque que l'indice de Jaccard pour les quadrigrammes multiplie la chance pour une paire d'être parallèle
ou semi-parallèle par 21. On observe aussi que l'indice de Jaccard sur les trigrammes est quant à lui négatif, multipliant les chances pour
une paire par 0,70, lui ôtant donc des chances d'être parallèle.

\begin{figure}[H]
    \centering
    \includegraphics[width=\textwidth]{3feature_importance.png}
    \caption{Importance de 3 variables utilisées par un modèle entraîné spécifiquement.}
    \label{3feature_importance}
\end{figure}

On remarque qu'un modèle entraîné uniquement sur les variables de score SBERT, Jaccard trigrammes et quadrigrammes
présente des scores similaires bien que moins élevés (voir figure \ref{3feature_importance}). Ici on remarque par exemple
que le trigramme reste pondéré négativement, tandis que les quadrigrammes restent la variable avec le plus de poids dans la classification.
Cependant nous pouvons aussi remarquer que l'absence des autres caractéristiques réduit considérablement le poids donné aux variables présentes.

\begin{figure}[H]
    \centering
    \includegraphics[width=\textwidth]{feature_importance_multiclass.png}
    \caption{Importance des variables utilisées par le modèle multiclasse.}
    \label{feature_importance_mc}
\end{figure}

Pour le modèle multiclasse (voir \ref{feature_importance_mc}), ces 2 vecteurs de coefficients montrent des profils distincts.
La classe parallèle est principalement pilotée
par les jaccard\_trigrams (2,20), qui constituent le prédicteur le plus fort. Les paires strictement parallèles
partagent des séquences de caractères de longueur 3 bien plus fréquentes que les autres classes.
La classe non-parallèle se distingue par des coefficients négatifs pour le sbert\_score (-0,80), le jaccard\_quadrigrams (-2,00)
et le phonetic\_jaccard (-1,00), confirmant que l'absence de similarité sémantique, lexicale et phonétique est le
marqueur principal du non-parallélisme. Le length\_score positif pour cette classe indique que les paires non-parallèles
ont tendance à présenter des écarts de longueur plus marqués.

\begin{figure}[H]
    \centering
    \includegraphics[width=\textwidth]{feature_importance_mc_or.png}
    \caption{Rapports de cotes des caractéristiques dans le modèle multiclasse.}
    \label{feature_importance_mc_or}
\end{figure}

Pour le modèle multiclasse, on observera les rapports de cotes par classe. La figure \ref{feature_importance_mc_or} présente les rapports de cotes pour
les classes non-parallèles et parallèles. On peut remarquer que pour qu'une paire soit considérée non-parallèle le score de longueur augmente grandement
ses chances là où les autres scores diminuent ses chances, surtout l'indice de Jaccard sur les quadrigrammes, le score SBERT et l'indice de Jaccard sur
la représentation phonétique. On remarque que pour qu'une paire soit considérée comme parallèle, le score de longueur devient très discriminant (ayant un 
rapport de cotes proche de 0,1) et l'indice de
Jaccard sur les trigrammes multiplie les chances par 9, augmentant grandement les chances pour une paire d'être parallèle si ces 2 caractéristiques présentent
un haut score.

\begin{figure}[H]
    \centering
    \includegraphics[width=0.65\textwidth]{feature_importance_multiclass_semi.png}
    \caption{Importance des variables utilisées par le modèle multiclasse pour la classe semi-parallèle.}
    \label{feature_importance_mc_semi}
\end{figure}

La classe semi-parallèle (voir \ref{feature_importance_mc_semi}) est en revanche
dominée par le jaccard\_quadrigrams (2,70) et le length\_score (1,70), suggérant que les
semi-parallèles se caractérisent par des similitudes lexicales combinées à des écarts de longueur significatifs.

Globalement, ces résultats confirment que le modèle apprend des distinctions
cohérentes avec le sens linguistique. Les parallèles stricts se caractérisent par des similitudes lexicales courtes et une similarité
sémantique modérée, les semi-parallèles par des similitudes lexicales longues mais des écarts de longueur, et les non-parallèles par
l'absence globale de similarité.

\begin{figure}[H]
    \centering
    \includegraphics[width=0.65\textwidth]{feature_importance_mc_or_semi.png}
    \caption{Rapports de cotes des caractéristiques pour la classe semi-parallèle.}
    \label{feature_importance_mc_or_semi}
\end{figure}

Dans le cas de la classe semi-parallèle on remarque que les rapports de cotes (voir figure \ref{feature_importance_mc_or_semi}) font pratiquement tous
augmenter les chances. Cependant il est intéressant de noter que le score de longueur vient, comme pour la classe non-parallèle, augmenter les chances
pour une paire de faire partie de la classe. L'indice de Jaccard sur les quadrigrammes présente la plus grande différence comparé aux rapports des autres
classes. En effet, celui-ci est extrêmement représentatif de la classe semi-parallèle au sein du modèle multiclasse comparé aux autres caractéristiques.
Cet indice est donc l'indice le plus discriminant pour qu'une paire soit assignée à la classe semi-parallèle.

\newpage

\subsection{Évaluation du modèle binaire}

\begin{table}[H]
\centering
\begin{tabular}{lSSS}
\hline
\textbf{Classe} & \textbf{Précision} & \textbf{Rappel} & \textbf{Score F1} \\
\hline
non-parallèle & 0,63 & 0,90 & 0,74 \\
parallèle & 0,93 & 0,74 & 0,82 \\
moyenne & 0,78 & 0,82 & 0,78 \\
moyenne pondérée & 0,83 & 0,79 & 0,79 \\
\hline
\end{tabular}
\caption{Rappel, précision et score F1 du modèle binaire.}
\label{binary_rpf1}
\end{table}

La moyenne pondérée est calculée en faisant la moyenne des métriques de chaque classe, pondérée par le nombre d'échantillons
de chaque classe. Concrètement, chaque classe contribue proportionnellement à sa fréquence dans le jeu de test, donc les classes
plus nombreuses ont plus d'influence sur le score final. Contrairement à la moyenne macro qui donne le même poids à chaque classe,
la moyenne pondérée reflète mieux la performance globale sur l'ensemble des données. La moyenne pondérée n'est pas égale à la moyenne
car les paires semi-parallèles sont aussi considérées dans la pondération (en l'occurrence du côté des parallèles).

Le modèle de régression logistique binaire atteint une précision de 83~\% en moyenne pondérée sur l'ensemble de test,
composé de 687 paires de phrases. La précision atteint 93~\% pour les paires parallèles ou semi-parallèles, contre 63~\%
pour les paires non-parallèles, ce qui indique que le modèle est plus fiable pour identifier les correspondances que pour
rejeter les correspondances. Le rappel est de 74~\% pour la classe parallèle et 90~\% pour les non-parallèles, montrant une
légère tendance à manquer quelques vraies non-correspondances. Le F1-score global pondéré de 79~\% confirme une performance robuste,
avec un F1 de 82~\% pour la classe parallèle contre 74~\% pour les non-parallèles.

\begin{figure}[H]
    \centering
    \includegraphics[width=0.75\textwidth]{confusion_binary.png}
    \caption{Matrice de confusion du modèle binaire.}
    \label{confusion_binary}
\end{figure}

La matrice de confusion (\ref{confusion_binary}) du modèle binaire (basée sur la partie d'évaluation de 687 paires)
montre des résultats globalement satisfaisants. Sur les 229 paires
réellement non-parallèles du jeu de test, 205 sont correctement classées, soit un taux de vrais négatifs de
89,5~\%. Les 24 faux positifs restants correspondent à des paires non-parallèles incorrectement identifiées
comme parallèles. Sur les 458 paires réellement parallèles ou semi-parallèles, 337 sont
correctement identifiées, soit un taux de vrais positifs de 73,6~\%, tandis que 121 sont incorrectement
rejetées comme non-parallèles. Ce taux de faux négatifs plus élevé s'explique en partie par la présence de
paires semi-parallèles dans la classe positive. Leurs scores sémantiques et de surface étant intermédiaires,
elles sont naturellement plus difficiles à distinguer des paires non-parallèles. Ce
déséquilibre entre le taux de vrais négatifs et le taux de vrais positifs indique que le modèle est légèrement
conservateur, préférant rejeter des paires douteuses plutôt que d'introduire du bruit dans les données
extraites. Ce comportement est acceptable dans notre contexte, où nous préférons la qualité à la quantité.

\begin{figure}[H]
    \centering
    \includegraphics[width=\textwidth]{roc_pr_binary.png}
    \caption{Courbes AUC-ROC et précision rappel du modèle binaire.}
    \label{roc_pr_binary}
\end{figure}

La figure \ref{roc_pr_binary} présente les \gls{courbes AUC-ROC} et précision-rappel du modèle binaire, qui confirment de bonnes
performances globales.
L'AUC-ROC de 0,851 indique qu'il y a 85,1~\% de chance que le modèle donne un score plus élevé
à des phrases réellement parallèles qu'à des phrases réellement non-parallèles. La courbe ROC
présente une montée initiale rapide, signifiant que le modèle atteint un taux de vrais positifs élevé à
faible coût en faux positifs, ce qui est particulièrement souhaitable. La courbe précision-rappel, avec un score de précision
moyenne de 0,915, est encore plus favorable. Elle montre que pour des valeurs de rappel allant jusqu'à
environ 0,5, la précision se maintient proche de 1,0, prouvant que le modèle peut être utilisé en mode
conservateur pour extraire un sous-ensemble de paires de très haute qualité avec très peu de bruit.
La dégradation de la précision n'intervient qu'au-delà d'un rappel d'environ 0,5, lorsque le modèle
commence à récupérer les paires les plus ambiguës. Ces résultats suggèrent que l'ajustement du seuil de
décision constitue un levier important pour adapter le comportement du classifieur selon que l'on privilégie
la précision ou le rappel.

L'évaluation du modèle binaire repose sur un rapport de classification incluant précision, rappel
et score F1 pour chaque classe, ainsi que sur une matrice de confusion permettant d'identifier les types
d'erreurs les plus fréquentes. La courbe ROC et la courbe précision-rappel complètent cette évaluation en
quantifiant la capacité du modèle à séparer les deux classes à différents seuils de décision. Ces deux
courbes sont particulièrement pertinentes dans notre contexte, où le coût d'un faux positif (une paire
non-parallèle incorrectement retenue) et celui d'un faux négatif (une paire parallèle incorrectement
rejetée) ne sont pas nécessairement symétriques selon l'usage visé. L'analyse des coefficients du modèle
permet enfin d'identifier quelles variables contribuent le plus à la décision de classification.

\subsection{Évaluation du modèle multiclasse}

\begin{table}[H]
\centering
\begin{tabular}{lSSS}
\hline
\textbf{Classe} & \textbf{Précision} & \textbf{Rappel} & \textbf{Score F1} \\
\hline
Non-parallèle & 0,68 & 0,77 & 0,72 \\
Semi-parallèle & 0,77 & 0,52 & 0,62 \\
Parallèle & 0,77 & 0,93 & 0,84 \\
Moyenne & 0,74 & 0,74 & 0,73 \\
Moyenne pondérée & 0,74 & 0,74 & 0,73 \\
\hline
\end{tabular}
\caption{Rappel, précision et score F1 du modèle multiclasse.}
\label{multiclass_rpf1}
\end{table}

La moyenne pondérée est calculée comme précédemment. On remarque dans ce cas que la moyenne pondérée est égale à la moyenne.
Cela est dû au fait que dans notre dataset il y a autant de représentants pour chacune des classes.

Le modèle multiclasse, qui distingue les trois catégories, affiche un score F1 moyen pondéré de 73~\%,
inférieur au modèle binaire, ce qui est cohérent avec la complexité accrue de la tâche. Les performances varient considérablement
selon les classes~: la classe parallèle est bien identifiée avec un F1 de 84~\% et un rappel de 93~\%, tandis que la classe
semi-parallèle présente les résultats les plus faibles avec un F1 de 62~\% et un rappel de seulement 52~\%. Ce faible rappel pour
les semi-parallèles signifie que près de la moitié de ces paires sont mal classées, soit comme parallèles, soit comme non-parallèles.
La classe non-parallèle obtient un F1 de 72~\%, comparable au résultat binaire pour la même catégorie. Cette disparité souligne
la difficulté à distinguer les paires semi-parallèles des deux autres classes.

\begin{figure}[H]
    \centering
    \includegraphics[width=0.75\textwidth]{confusion_multiclass.png}
    \caption{Matrice de confusion du modèle multiclasse.}
    \label{confusion_multiclass}
\end{figure}

La matrice de confusion \ref{confusion_multiclass} du modèle multiclasse révèle, elle,
une asymétrie dans les performances selon les
classes. Les paires parallèles sont correctement classées dans 212 cas sur 229, soit un taux de 92,6~\%,
ce qui confirme que le modèle identifie bien les correspondances exactes. Les paires non-parallèles
obtiennent un taux de classification correcte de 76,8~\% (176 sur 229), les erreurs se répartissant presque égalitairement
entre des confusions avec la classe parallèle (24 cas) et la classe semi-parallèle
(26 cas). Les paires semi-parallèles constituent sans surprise la classe la plus difficile à identifier
correctement~: seulement 119 sur 229 sont bien classées, soit 51,9~\%. Les erreurs se concentrent massivement
vers la classe non-parallèle (77 cas). Cette confusion reste acceptable dans notre cas pour la même raison que lorsque
des paires semi-parallèles étaient classées comme non-parallèles par le modèle binaire. Les 38 confusions entre
semi-parallèle et parallèle correspondent probablement à des cas où la troncature était mineure et n'a
pas suffisamment dégradé les variables pour distinguer les deux catégories. Ces résultats montrent que la
frontière entre paires semi-parallèles et non-parallèles est la plus difficile à établir, ce qui est cohérent
avec la nature artificielle des semi-parallèles présentes dans nos sous-ensembles.

\begin{table}[H]
\centering
\begin{tabular}{l p{0.35\linewidth} p{0.35\linewidth}}
\hline
\textbf{Indice} & \textbf{FR} & \textbf{LR} \\
\hline
0 &
Considérant que la reconnaissance de la dignité inhérente à tous les membres de la famille humaine et de leurs droits égaux et inaliénables constitue le fondement de la liberté, de la justice et de la paix dans le monde, &
Ella, allora si e astizzatu, e nun vulia più entra in casa.Lu Babbu allora esciutu fora e à cuminciatu a prigallu d'entra in casa. \\
1 &
Les médecins l'ont montré~: cette discrimination peut engendrer des problèmes psychologiques, voire des maladies organiques, comme des cancers de la face - le visage étant le siège de l'identité.&
A la fin, estent dintrat en el mèmes, diguèt : quantes i a dins l’ostal de mon paire de vailets que an de pan mai lo que lor’n cal, et ieu aicí aganissi.\\
2 &
C'est là l'un des éléments essentiels qu'oublient ceux qui ne voient pas l'intérêt de sauver le tariana (Amazonie), le thouktche (Sibérie) ou le poitevin saintongeais~: une langue ne permet pas seulement de transmettre une information ou de trouver du travail.&
Eccu unu di l’elementi essenziali scurdatu da quelli ch’ùn vedenu l’interessu di salvà u tariana (Amazonia)\\
3 &
L'enseignement technique et professionnel doit être généralisé; l'accès aux études supérieures doit être ouvert en pleine égalité à tous en fonction de leur mérite.&
De Fàch-ùn Beruefsùnterrìcht muess àllgemein zügangli sìnn. Àlli mìen de salwe offene Zügàng zuem Hochschuelùnterrìcht hàn\\
4&
Un homme avait deux fils. Dont le plus jeune dit à son père~: mon père, [...]&
Un omë avié dous éfans. Lou pus jhouiné diguet à soun pero~: Moun pero, [...]\\
5&
Pour quoi, comme les bonnes choses peuvent toujours être utiles, on doit les écouter avec attention, quel que soit le narrateur.&
Drùm muess mr de guete Gedànke ùffmerriksàm züheere, egàl wër se verzehlt, dann sie kénne ìmmer nìtzli sìnn.\\
\hline
\end{tabular}
\caption{Phrases et leur indice pour comprendre le tableau \ref{good_bad_classes}.}
\label{sentences_for_gbc}
\end{table}

\begin{table}[H]
\centering
\begin{tabular}{lScc}
\hline
\textbf{Indice} & \textbf{Confiance} & \textbf{Classe réelle} & \textbf{Prédiction} \\
\hline
0 & \qty{91,26}{\percent} & non-parallèle & non-parallèle \\
1 & \qty{90,20}{\percent} & non-parallèle & parallèle \\
2 & \qty{91,62}{\percent} & semi-parallèle & semi-parallèle \\
3 & \qty{83,53}{\percent} & semi-parallèle & non-parallèle \\
4 & \qty{99,89}{\percent} & parallèle & parallèle \\
5 & \qty{62,16}{\percent} & parallèle & non-parallèle \\
\hline
\end{tabular}
\caption{Exemples de phrases bien et mal classées}
\label{good_bad_classes}
\end{table}

Le tableau \ref{good_bad_classes} représente la classification et le score de confiance des phrases présentées dans le tableau \ref{sentences_for_gbc}.
La phrase non-parallèle la mieux classée oppose deux textes qui ne partagent aucune similarité sémantique.
Le modèle la prédit non-parallèle avec une certitude de 91,26~\%, identifiant correctement 
l'absence de lien entre les phrases. La moins bien classée, quant à elle, est prédite comme parallèle avec une certitude de 90~\%.
Cette classification doit provenir de la longueur similaire des phrases et de la structure (de par la ponctuation), elle aussi similaire.
Ce faux positif est très mauvais dans notre cas car il démontre une incapacité du modèle à classer certaines phrases non-parallèles correctement.

La phrase semi-parallèle la mieux classée avec une certitude de 92~\% est une traduction presque complète de
la phrase française. Le modèle identifie correctement ce semi-parallélisme dû à la similarité sémantique et structurelle de la phrase.
La phrase la moins bien classée est prédite comme non-parallèle avec une certitude de
84~\%, confondant une traduction partielle avec une non-correspondance. C'est un faux négatif probablement lié à la barrière orthographique
de l'alsacien (dû aux racines germaniques de la langue) et à la faible similarité structurelle entre les deux phrases.

La phrase parallèle la mieux classée (score de confiance de 99,89~\%) est une traduction occitane qui suit fidèlement la structure française.
Le modèle identifie cette correspondance structurelle avec une certitude quasi absolue.
La moins bien classée est prédite comme étant non-parallèle avec une certitude de 62~\%. C'est un faux négatif~: il s'agit bien d'une traduction
parallèle, mais la barrière orthographique de l'alsacien (caractères spécifiques comme ù, ì, è) réduit la similarité de caractères à zéro.
Nous pouvons remarquer que c'est le score de certitude le plus bas que nous avons pu remarquer pour un faux négatif.

Nous remarquons que les langues ayant une orthographe très distincte du français (comme l'alsacien) produisent
souvent des faux négatifs (prédictions non-parallèles au lieu de semi-parallèles ou parallèles dans les contextes où les phrases devraient l'être)~:
la similarité phonétique ou structurelle est trop faible pour compenser la
similarité sémantique basse. En effet, le modèle prend souvent les semi-parallélismes pour des parallélismes, ou
non-parallélismes, selon le score des différentes variables. Ces limites révèlent que la régression logistique, bien que performante, reste
sensible à la qualité des attributs linguistiques et ne capture pas toujours les subtilités linguistiques permettant une classification correcte
des phrases.

\section{Validation du modèle sur un nouveau corpus}

Afin d’évaluer la qualité du modèle obtenu, nous procédons à une validation externe en utilisant le corpus
de \citet{tsolakis_2024} issu de données Wikipédia. Il contient des paires de phrases alignées en corse, occitan,
alsacien et français. Il présente deux caractéristiques importantes pour notre évaluation. D'une part, il est
structurellement distinct du corpus d'entraînement, reposant sur des textes encyclopédiques hétérogènes
plutôt que sur des textes narratifs ou littéraires. D'autre part, il couvre les mêmes paires de langues que
notre corpus d'entraînement, ce qui permet une évaluation directe de la généralisation du modèle.

L’application du modèle sur ce corpus permet de vérifier~:
\begin{itemize}
    \item la capacité du classifieur à identifier correctement des phrases parallèles dans un contexte
    différent
    \item l’absence de surapprentissage spécifique au corpus LIMSI et à ses textes majoritairement occitans
    \item la robustesse du modèle face aux différences structurelles et thématiques.
\end{itemize}

La phase de validation constitue l'étape finale de la chaîne de traitement, permettant d'évaluer la capacité du modèle
à identifier des fragments parallèles dans un corpus comparable réel, structurellement et thématiquement
distinct du corpus d'entraînement. Elle repose sur deux programmes distincts dont la séparation est
intentionnelle~: un script de préparation des données d'une part, et un notebook d'inférence et d'analyse
d'autre part.

\newpage

\subsection{Préparation du corpus de validation} \label{prepa_corpus_v}

Les fichiers textes bruts issus des dossiers de validation sont d'abord transformés en tableaux
structurés de paires de phrases exploitables par le classifieur. Ce découpage en deux étapes présente un
avantage important~: le nettoyage et la segmentation, potentiellement coûteux en temps de calcul
sur de grands corpus, ne sont effectués qu'une seule fois. Les relances ultérieures de validation
repartent directement des fichiers CSV ainsi produits, sans nécessiter de retraitement des textes sources.

Les fichiers français et régionaux sont appariés par préfixe de nom
de fichier commun, garantissant que les combinaisons générées restent intra-documentaires, c'est-à-dire
que seules des phrases issues d'articles traitant du même sujet sont mises en regard. Cette contrainte
est essentielle dans le cadre d'un corpus comparable. Comparer des phrases issues de documents sans
lien thématique produirait une proportion de paires non-parallèles trop importante et rendrait l'extraction de
fragments pertinents impossible.

Les textes bruts présentent une structure caractéristique des extraits encyclopédiques.
Une chaîne de nettoyage en deux étapes est donc appliquée avant toute segmentation. La première étape
supprime les lignes non informatives (titres courts de quatre mots ou moins sans ponctuation finale,
éléments de liste réduits à un nom propre isolé, et lignes constituées d'une simple puce ou d'un
tiret sans contenu réel). Ces éléments, s'ils étaient conservés, produiraient des segments trop courts ou
trop vides d'information pour être intéressants à comparer. La seconde étape normalise les espaces multiples
et les tirets longs afin d'homogénéiser la ponctuation.

La segmentation repose sur un découpage aux ponctuations fortes, c'est-à-dire le point, le
point d'exclamation et le point d'interrogation, suivies d'une majuscule. Afin d'éviter les coupures
erronées sur des abréviations ou des nombres, un mécanisme de protection par
substitution temporaire est appliqué avant le découpage. Les abréviations identifiées sont remplacées
par des marqueurs neutres, les nombres décimaux sont protégés de la même manière. Ces substitutions
sont annulées après la segmentation. Les phrases de moins de cinq tokens sont ensuite écartées, car elles
ne portent pas suffisamment d'information (sémantique ou lexicale) pour alimenter les variables d'apprentissage.

Une fois les phrases extraites, on génère l'ensemble des combinaisons entre phrases
françaises et phrases régionales au sein de chaque paire de documents. Le produit, bien que
potentiellement volumineux pour des documents longs, est la stratégie naturelle pour un corpus comparable
non aligné. En l'absence d'alignement préalable, toutes les combinaisons doivent être évaluées par le
classifieur afin d'identifier celles qui présentent un parallélisme suffisant. Le résultat est
exporté sous forme de CSV structuré, avec une ligne par paire de phrases et des colonnes indiquant la
variété linguistique, le code langue, la paire de documents d'origine, la phrase française et la phrase
régionale.

\begin{figure}[H]
    \centering
    \includegraphics[width=\textwidth]{extract_sentence_pairs_comparable.png}
    \caption{Extrait du CSV des paires de phrases.}
    \label{extract_sentence_pairs_comparable}
\end{figure}

Voici ce que donne l'appairage des phrases comparables sous forme de CSV (voir figure \ref{extract_sentence_pairs_comparable}).
Chaque phrase en français est appariée, dans cet exemple, à une phrase en alsacien pour que le programme de validation
infère ensuite si une phrase est parallèle, semi-parallèle ou non-parallèle et la classer en fonction.
Ces données sont triées pour chaque paire de documents fr-lr (lr pour langue régionale) et réalisées pour chaque paire langagière (voir tableau \ref{pairs}).

\begin{table}[H]
\centering
\begin{tabular}{lS}
\hline
\textbf{Dialectes} & \textbf{Nombre de paires} \\
\hline
fr-co & 1\ 997 \\
fr-gsw & 572\ 372 \\
fr-gsw\_mulhousien & 323\ 733 \\
fr-gsw\_strasbourgeois & 2\ 781 \\
fr-oc & 1\ 040\ 936 \\
\hline
\end{tabular}
\\
\smallskip
\small{(fr: français, co: corse, gsw: alsacien, oc: occitan)}
\smallskip
\caption{Nombre de paires de phrases par paires langagières.}
\label{pairs}
\end{table}

\subsection{Infèrence et analyse}

Pour la validation, nous chargeons les CSV produits par le script de préparation. Nous calculons ensuite les variables
d'apprentissage sur l'ensemble des paires, puis appliquons les deux modèles entraînés (binaire et multiclasse). 
Nous calculons pour chaque paire les mêmes variables que celles que nous avons calculées pour le corpus d'entraînement.
Cette étape est indispensable pour garantir
que les données présentées au classifieur soient comparables à celles sur lesquelles il
a été entraîné. Les variables sont ensuite normalisées à l'aide du scaler ajusté lors de l'entraînement.

L'inférence produit pour chaque paire une prédiction binaire opposant les paires parallèles ou semi-parallèles
aux paires non-parallèles, accompagnée d'un score de confiance inféré sur la
probabilité que le modèle attribue à la classe prédite pour chaque paire, ainsi qu'une prédiction multiclasse
distinguant les trois catégories. Le score de confiance du modèle binaire est utilisé comme critère de
filtrage. Seules les paires dont la confiance dépasse un seuil paramétrable, fixé par défaut à 0,75, sont
retenues comme candidates. Ce seuil constitue le principal levier d'ajustement du comportement du système.
Celui-ci a été choisi en se basant sur l'hypothèse qu'un tel seuil serait intéressant car il ne prend en
compte que les paires ayant un score supérieur à 75~\% de confiance, ce qui, théoriquement, devrait retirer
la vaste majorité des paires non-parallèles au vu des résultats provenant de nos variables vues plus haut.

L'analyse des résultats est conduite à deux niveaux de granularité. Au niveau des variétés, les statistiques
permettent de comparer les performances du modèle entre l'alsacien, le corse et l'occitan et entre les
variétés dialectales potentiellement présentes pour ces 3 langues. Cette séparation est pertinente dans
notre contexte, car les variétés mulhousienne et strasbourgeoise de l'alsacien
(annotées comme telles par la page Wikipédia le précisant dans son code de langue) peuvent présenter des graphies
sensiblement différentes de l'alsacien "par défaut" (pages du Wikipédia alémanique notées comme étant en "alsacien")
de Wikipédia, ce qui est susceptible d'affecter les variables de surface
tels que les indices de Jaccard. Au niveau des langues, la comparaison permet d'évaluer dans quelle
mesure le modèle, entraîné généralise correctement à des langues
structurellement différentes les unes des autres.

Enfin, les paires candidates retenues sont exportées sous forme de CSV, à la fois dans un fichier global
et dans des fichiers individuels par variété. Ces fichiers constituent la sortie opérationnelle du système~:
ils contiennent les fragments parallèles et semi-parallèles identifiés automatiquement dans le corpus
comparable, accompagnés de l'ensemble des scores calculés, et peuvent être utilisés directement pour enrichir
les ressources parallèles disponibles pour les langues régionales étudiées.

\subsection{Résultats de la validation} \label{results_validation}

\begin{figure}[H]
    \centering
    \includegraphics[width=\textwidth]{heatmap_sbert_validation.png}
    \\
    \smallskip
    \small{(classes~: 1~: parallèle, 2~: semi-parallèle)}
    \smallskip
    \caption{Score SBERT moyen par classe et dialecte.}
    \label{heatmap_sbert_validation}
\end{figure}

Cette carte de chaleur (\ref{heatmap_sbert_validation}) présente les scores SBERT moyens des paires candidates retenues lors de la validation,
ventilés par dialecte et par classe prédite par le modèle multiclasse. Les classes sont encodées
numériquement~: 1 correspond aux paires parallèles et 2 aux paires semi-parallèles. Contrairement à
la heatmap produite lors de l'évaluation sur le jeu de test, la classe 0 (non-parallèle) est absente
de cette figure, ce qui signifie qu'aucune paire non-parallèle n'a franchi le seuil de confiance pour
aucune des variétés. L'absence totale de la classe 0 ici suggère que le
modèle, appliqué au corpus Wikipédia, filtre efficacement les paires avec un lien sémantique très limité.

Les scores des classes parallèle et semi-parallèle sont globalement cohérents avec ceux observés lors de
l'évaluation sur le jeu de test, avec quelques différences notables. Le corse conserve les scores les plus
élevés pour les deux classes, à 0,57 et 0,58 respectivement, confirmant la tendance observée précédemment
et la meilleure captation de la proximité sémantique entre le français et cette langue romane par le modèle
SBERT. Le strasbourgeois présente des scores intermédiaires à 0,52 et 0,53, nettement supérieurs à ceux de
l'alsacien "par défaut" et du mulhousien qui se situent entre 0,44 et 0,47. Cette hiérarchie au sein des variétés
alsaciennes suggère que certains dialectes au sein des langues régionales (ici dans le cas de l'alsacien, le strasbourgeois)
produisent des paires candidates sémantiquement plus proches du français. L'occitan affiche un score parallèle de 0,50,
légèrement supérieur à celui de
l'alsacien standard, mais un score semi-parallèle de seulement 0,45, inférieur à celui du strasbourgeois.

On notera également que les scores des classes parallèle et semi-parallèle sont très proches pour
chaque dialecte, l'écart maximal observé étant de 0,05 pour le mulhousien. Cette proximité confirme
la difficulté structurelle à distinguer ces deux catégories par le seul score sémantique, déjà identifiée
lors de l'analyse des matrices de confusion et des distributions KDE.

\begin{figure}[H]
    \centering
    \includegraphics[width=0.85\textwidth]{Sbert_distribution_by_lp_oc_co.png}
    \caption{Distribution des scores SBERT moyen pour l'occitan et le corse.}
    \label{Sbert_distribution_by_lp_oc_co}
\end{figure}

Les figures \ref{Sbert_distribution_by_lp_oc_co} présentent deux histogrammes sur la distribution des scores SBERT pour le corse et l'occitan,
superposant en gris l'ensemble des comparaisons effectuées et en vert les seules paires candidates retenues.
La ligne pointillée verticale indique le score SBERT moyen des paires candidates.

\begin{figure}[H]
    \centering
    \includegraphics[width=0.85\textwidth]{Sbert_distribution_by_lp_gsw.png}
    \caption{Distribution des scores SBERT moyen pour les dialectes alsaciens.}
    \label{Sbert_distribution_by_lp_gsw}
\end{figure}

Les figures \ref{Sbert_distribution_by_lp_gsw} présentent trois histogrammes, un par dialecte alsacien, pour montrer la
distribution des scores SBERT comme sur la figure \ref{Sbert_distribution_by_lp_oc_co}.

La première observation transversale à toutes les langues est la concentration massive des comparaisons
totales dans les valeurs très basses, entre 0 et 0,2. Ce phénomène est inhérent à notre stratégie.
La très grande majorité des combinaisons de phrases générées au sein d'un même document
comparable sont sémantiquement éloignées, ce qui confirme que le corpus utilisé est bien un corpus
comparable et non un corpus parallèle. Le volume de comparaisons varie considérablement selon les langues,
les histogrammes de l'alsacien "par défaut" et de l'occitan atteignant respectivement environ 140 000 et
260 000 comparaisons, montrant leur présence importante au sein du corpus comparable.

Les paires candidates, représentées en vert, apparaissent systématiquement dans les valeurs intermédiaires
à hautes du score SBERT, avec des moyennes comprises entre 0,42 et 0,54 selon les langues. Pour le corse
et le strasbourgeois, dont les corpus sont plus petits, les candidates sont peu nombreuses mais bien
concentrées au-delà de 0,4, indiquant que le filtre a correctement ciblé les paires sémantiquement
proches. Pour l'occitan et les variétés alsaciennes, la distribution des candidates est plus étalée, ce qui
s'explique par le volume plus important de comparaisons disponibles et donc une diversité plus importante de
paires potentiellement ambiguës atteignant le seuil de confiance.

\begin{figure}[H]
    \centering
    \includegraphics[width=\textwidth]{trust_threshold.png}
    \caption{Seuil de confiance par sbert moyen et nombre de candidats.}
    \label{trust_threshold}
\end{figure}

Le graphique à double axe \ref{trust_threshold} permet d'observer simultanément l'évolution du nombre de paires candidates et
du score SBERT moyen de ces candidates en fonction du seuil de confiance appliqué, le seuil actuel de 0,75
étant matérialisé par la ligne pointillée rouge. Le nombre de paires candidates décroît de manière monotone et
régulière à mesure que le seuil augmente,
passant d'environ 275\,000 paires à un seuil de 0,50 à environ 25\,000 paires à un seuil de 0,95. La
décroissance est approximativement linéaire entre 0,50 et 0,75, puis ralentit légèrement au-delà,
indiquant qu'un noyau de paires bénéficie d'une confiance très élevée du modèle. Au seuil actuel de 0,75,
environ 100\,000 paires sont retenues (toutes langues confondues).

Le score SBERT moyen des candidates évolue en sens inverse, progressant d'environ 0,35 à un seuil de 0,50 jusqu'à
environ 0,70 à un seuil de 0,95. Cette progression confirme que l'élévation du seuil de confiance permet
effectivement de sélectionner des paires présentant une proximité sémantique plus élevée. On observe
cependant que la progression du score SBERT moyen s'accélère au-delà de 0,75, passant de ~0,5
à ~0,7 entre 0,75 et 0,95 de seuil de confiance, alors que le gain entre 0,50 et 0,75 est plus modeste (de ~0,4 à ~0,5).
Ce point autour de 0,75 constitue une information utile pour le réglage du seuil. Un seuil dans cette
zone permet de réduire le volume de paires candidates tout en garantissant une bonne qualité
sémantique.

\begin{table}[H]
\centering
\begin{tabular}{lSSSSSS}
\hline
\textbf{Dialectes} & \textbf{Candidates} & \textbf{Parallèles} & \textbf{Semi} & \textbf{Extraction} & \textbf{Conf. moy.} & \textbf{SBERT moy.} \\
\hline
fr-co & 124 & 35 & 89 & \qty{6,21}{\percent} & \qty{86,30}{\percent} & 0,5809 \\
fr-gsw & 13\ 019 & 3\ 558 & 9\ 461 & \qty{2,27}{\percent} & \qty{85,90}{\percent} & 0,4469 \\
fr-mul & 8\ 445 & 2\ 228 & 6\ 217 & \qty{2,61}{\percent} & \qty{85,86}{\percent} & 0,4618 \\
fr-str & 114 & 24 & 90 & \qty{4,10}{\percent} & \qty{85,98}{\percent} & 0,5294 \\
fr-oc & 78\ 195 & 36\ 300 & 41\ 895 & \qty{7,51}{\percent} & \qty{87,49}{\percent} & 0,4752 \\
\hline
\end{tabular}
\\
\smallskip
\small{(fr: français, co: corse, gsw: alsacien, oc: occitan, mul: mulhousien, str: strasbourgeois)\\}
\small{Extraction~: pourcentage de phrases extraites sur l'entiereté du corpus}
\smallskip
\caption{Performances du modèle sur la validation par dialectes.}
\label{perfs_dialect}
\end{table}

Le tableau \ref{perfs_dialect} présente les statistiques sur les phrases extraites par le modèle évaluées par
dialecte sur l'ensemble de validation. Le
dialecte fr-oc (occitan) constitue la classe la plus représentée avec 78\,195 paires
candidates extraites, dont 36\,300 étiquetées comme parallèles, ainsi que le
meilleur score de confiance moyen 87,49~\% et le score SBERT le plus
élevé 0,4752. À l'opposé, le dialecte fr-co (corse), bien que présentant la
confiance moyenne la plus faible (86,30~\%), affiche le score SBERT le
plus élevé (0,5809), suggérant une meilleure adéquation sémantique entre les
phrases des paires fr-co candidates, malgré un nombre de candidates bien plus bas
(124 candidates dont 35 étiquetées parallèles). 

Sur le plan qualitatif, on remarque
que le pourcentage de phrases parallèles candidates varie considérablement, allant de
2,27~\% pour fr-gsw à 7,51~\% pour fr-oc, ce qui reflète l'hétérogénéité de la
complexité d'annotation entre les dialectes. Les dialectes fr-gsw (alsacien) et
fr-mul (mulhousien) affichent des scores SBERT modestes (0,4469 et 0,4618
respectivement) malgré des volumes de données plutôt importants, ce qui pourrait
indiquer une difficulté accrue du modèle à capturer les spécificités lexicales
et syntaxiques de ces langues germaniques. Globalement, les scores de confiance
moyens restent élevés et homogènes (85,86~\% à 87,49~\%), ce qui témoigne de
la robustesse générale du modèle, tandis que les variations du score SBERT
mettent en lumière les défis restants pour certains dialectes moins bien
représentés ou très différents du français en termes d'orthographe et de structure.

\subsection{Analyse globale des résultats}

\begin{figure}[H]
    \centering
    \includegraphics[width=\textwidth]{top_candidates_oc.png}
    \caption{Top 5 des paires candidates occitanes.}
    \label{top_candidates_oc}
\end{figure}

On remarque que les candidates occitanes (présentées dans la figure \ref{top_candidates_oc}) ont le plus haut score de confiance pour des
phrases considérées comme parallèles par le modèle binaire et multiclasse.
Cette langue étant la plus proche phonétiquement et sémantiquement du français, cette observation est parfaitement cohérente.
Cependant il est à noter que les paires candidates occitanes doivent être triées après détection. En effet, la communauté Wikipédia occitan a écrit
la version en occitan dans l'article mais redonne toujours la version française en dessous (voir figure \ref{top_candidates_oc_prbl}). Certaines des phrases
extraites sont donc des paires français-français, ce qui pose un énorme problème pour nous et les paires doivent donc être triées.

\begin{figure}[H]
    \centering
    \includegraphics[width=0.75\textwidth]{top_candidates_oc_prbl.png}
    \caption{Présentation d'un article wikipédia français et occitan}
    \label{top_candidates_oc_prbl}
\end{figure}


\begin{figure}[H]
    \centering
    \includegraphics[width=\textwidth]{top_candidates_co.png}
    \caption{Top 5 des paires candidates corse.}
    \label{top_candidates_co}
\end{figure}

On remarque dans le cas du corse (présenté dans la figure \ref{top_candidates_co}) que le score SBERT est bien plus bas que chez les phrases en occitan.
De plus, les phrases retenues sont, en général, relativement courtes et ne présentent que peu de mots en commun. La dernière phrase
présentée comme parallèle
présente d'ailleurs un problème. Son score de confiance est très haut, mais la phrase semble être considérée comme parallèle uniquement
due à la présence du nom d'une {\oe}uvre "Oral Fixation, Vol.2 (2005)" qui est équivalent entre le français et le corse.
On peut aussi remarquer que la troisième phrase, bien que très intéressante, présente un
score de confiance multiclasse très bas. Cependant la paire est en effet semi-parallèle et reste
intéressante à considérer comme une paire candidate.

\begin{figure}[H]
    \centering
    \includegraphics[width=\textwidth]{top_candidates_gsw.png}
    \caption{Top 5 des paires candidates alsaciennes.}
    \label{top_candidates_gsw}
\end{figure}

Dans la figure \ref{top_candidates_gsw}, nous pouvons voir les candidates alsaciennes. Ces phrases sont très courtes et sont en général
des listes de termes équivalents entre français et alsacien. Cette structure notable pose un problème dans notre cas et une
utilisation du programme pour l'entraînement de modèles de traduction devrait faire attention à ce genre de cas. Utiliser des listes
d'oeuvres ou d'autres entités nommées peut poser d'énormes problèmes durant un entraînement, car ceux-ci sont un cas particulier
où les noms utilisés sont souvent les mêmes entre les langues. Nous pouvons aussi faire la remarque que ces candidates possèdent certainement
cette structure car le modèle a du mal à trouver des phrases similaires entre le français (langue romane) et l'alsacien (langue germanique).
Il semble donc garder les phrases ayant une forte similarité sémantique et phonétique provenant de la proximité orthographique des
termes (dans les exemples de la figure, des entités nommées).

\begin{figure}[H]
    \centering
    \includegraphics[width=\textwidth]{top_candidates_gswS.png}
    \caption{Top 5 des paires candidates strasbourgeoise.}
    \label{top_candidates_gswS}
\end{figure}

Dans la figure \ref{top_candidates_gswS}, nous pouvons remarquer les mêmes problèmes que dans le cadre des phrases alsaciennes "général".
Il nous semble donc légitime de dire que le modèle a des difficultés à trouver des paires candidates intéressantes entre français et alsacien.

\begin{figure}[H]
    \centering
    \includegraphics[width=\textwidth]{top_candidates_gswM.png}
    \caption{Top 5 des paires candidates mulhousienne.}
    \label{top_candidates_gswM}
\end{figure}

On remarque dans le cas de l'alsacien mulhousien (figure \ref{top_candidates_gswM}) que celui-ci possède les pires scores de confiance multiclasse.
Ces scores sont indicatifs de la difficulté pour le modèle multiclasse à catégoriser ses paires comme des paires candidates.
Comme les autres formes d'alsacien il semblerait que la structure même de la langue soit trop différente du français pour que
le modèle trouve des paires intéressantes qui ne sont pas pleines d'entités nommées équivalentes entre français et alsacien.

Pour parvenir aux paires de la figure \ref{top_candidates_gswM}, nous avons dû limiter légèrement quelles paires
deviendraient des candidates. En effet, on remarque que certaines
des phrases provenant du Wikipédia en mulhousien ne sont en fait pas traduites mais simplement des phrases en français.
Ainsi nous avons empêché les paires présentant un score SBERT parfait et les phrases équivalentes
de devenir des paires candidates (voir la figure \ref{code_candidates_limitation}). Dans l'exemple ci-dessus, nous avons laissé la 3\up{e}
phrase telle quelle pour présenter le problème. Cependant il est à noter que la première paire candidate était une paire
non traduite et que celle-ci a donc été filtrée. Ce manque de traduction arrive parfois sur Wikipédia et prouve
être un point très problématique dans le cadre du TAL en général. Cette ressource devient parfois fausse à cause d'un manque
de traduction d'une langue à une autre alors que l'article possède un lien interlangue et que celui-ci existe dans les deux langues
(en théorie).

\begin{figure}[H]
    \centering
    \includegraphics[width=\textwidth]{code_candidates_limitation.png}
    \caption{Code de limitation des candidates.}
    \label{code_candidates_limitation}
\end{figure}

Les visualisations produites dans \ref{results_validation} présentent une analyse quantitative de l'inférence
du modèle (voir \ref{prepa_corpus_v}). La distribution des scores SBERT par
dialecte permet d'observer si les paires candidates présentent des scores significativement supérieurs à
l'ensemble des comparaisons effectuées, ce qui constituerait une validation indirecte de la pertinence du
filtrage. La courbe d'effet du seuil de confiance permet, quant à elle, d'identifier le point à
partir duquel une élévation du seuil entraîne une réduction rapide du nombre de candidates sans gain
substantiel en qualité moyenne, informant ainsi le choix du seuil opérationnel.

L'analyse croisée des résultats quantitatifs produits permet de présenter plusieurs observations.
Premièrement, la distribution des scores SBERT confirme de manière cohérente la hiérarchie attendue entre les trois classes à chaque étape~:
entraînement, évaluation sur le jeu de test et validation sur Wikipédia. Les médianes de 0,17, 0,40 et 0,60 observées pour les
paires non-parallèles, semi-parallèles et parallèles sur le corpus d'entraînement se retrouvent approximativement dans les résultats de
validation, avec des scores moyens de paires candidates compris entre 0,42 et 0,54 selon les variétés. Cette stabilité suggère que le modèle SBERT,
bien qu'entraîné majoritairement sur des langues bien dotées, capte une partie du signal sémantique pertinent pour les langues
régionales étudiées.

Deuxièmement, la variable
la plus discriminante s'avère être l'indice de Jaccard sur les quadrigrammes de caractères, avec un coefficient d'environ 3,0 dans
le modèle binaire, devançant le score SBERT (1,4) et l'indice phonétique (1,0). Ce résultat est contre-intuitif au regard de la littérature,
qui tend à présenter les méthodes sémantiques comme centrales. Il s'explique cependant par la nature des langues étudiées. En effet, l'occitan et le
corse partagent avec le français une graphie suffisamment proche pour que les chevauchements de
séquences de quatre caractères constituent un signal fiable. Ce phénomène semble être moins marqué pour l'alsacien, langue germanique, ce qui
explique en partie les scores de discrimination plus faibles observés pour la paire fr-gsw, avec une marge entre paires parallèles et
non-parallèles de seulement 0,30 contre 0,51 pour le corse.

Troisièmement, la distribution quasi-parfaite des classes pour les paires
fr-oc (34,3~\% chacune) et fr-gsw (35,5~\% de non-parallèles), et les performances du
classifieur sur le corse dans la phase de validation 
restent fragiles en raison du très faible volume de ce corpus (22 fichiers, 9 004 tokens).
Ces observations plaident pour un effort d'acquisition de texte comparable fr-lr ciblé sur le corse et l'alsacien avant toute extension du système.
Enfin, la courbe d'effet
du seuil de confiance révèle qu'un seuil de 0,75 constitue un point d'équilibre raisonnable entre volume et qualité
sémantique des candidates. En dessous de
ce seuil, la quantité augmente rapidement mais la qualité sémantique moyenne stagne autour de 0,40 à 0,50. Au-delà, la progression en qualité
s'accélère mais le volume chute fortement, réduisant l'utilité pratique du système.